% Political Intrigue Campaign Framework −−− Fate's Edge (v0.3)
% Expanded with corrected Fenwood genealogy, timeline notes, and "Blood on the Belworth" adventure.
\documentclass[11pt]{article}
\usepackage[margin=1in]{geometry}
\usepackage{enumitem}
\usepackage{hyperref}
\usepackage{xcolor}
\usepackage{titlesec}
\usepackage{array}
\usepackage{setspace}
\usepackage{graphicx}
\usepackage{longtable}
\usepackage{lettrine}
ℎypersetup{colorlinks=true, linkcolor=black, urlcolor=blue, citecolor=black}
\setlist{noitemsep, topsep=4pt, leftmargin=1.2em}
\titleformat{\section}{\normalfontŁarge\bfseries}{}{0pt}{}
\titleformat{\subsection}{\normalfontłarge\bfseries}{}{0pt}{}
\titleformat{⫕section}{\normalfont\normalsize\bfseries}{}{0pt}{}

% Custom commands
≠wcommand{\timer}[2]{\textbf{#1:} [#2]}
≠wcommand{\position}[1]{\textit{Position:} #1}
≠wcommand{\dv}[1]{\textit{DV:} #1}

\begin{document}
\begin{center}
{ŁARGE \textbf{Political Intrigue Campaign Framework}}\\[3pt]
{łarge for \emph{Fate's Edge} −−− Version 0.3 (Expanded)}\\[2pt]
\end{center}
\vspace{0.25em}
ℎrule\vspace{0.75em}

\section*{Design Goals}
\begin{itemize}
  ℹtem \textbf{Why:} Enable long−form play focused on courts, city−states, guilds, and synods where choices echo across seasons.
  ℹtem \textbf{Mechanical Need:} Formalize social conflict and faction play using existing engines (Position, DV, SB, timers).
  ℹtem \textbf{Narrative Richness:} Turn every roll into a living thread−−−alliances, scandals, mandates, and crises.
  ℹtem \textbf{Player Appeal:} Provide levers (Favor, Leverage, Reputation, Strings) without heavy bookkeeping.
\end{itemize}
\emph{Principle:} This module adds procedures, not new math. Use core rules for rolls; this tells you what to prep, track, and tick.

\section*{Quickstart (2 minutes)}
\begin{enumerate}
  ℹtem Pick a \textbf{Seat of Power} (city, court, guildhall, synod).
  ℹtem Create \textbf{3−−5 Faction Sheets}. Give each faction an \textbf{Influence} timer [6], an \textbf{Exposure} timer [6], and a \textbf{Stability} timer [6−−8].
  ℹtem Place \textbf{Campaign Dials}: \emph{Mandate} and \emph{Crisis} at 2 (0−−6 each) unless fiction says otherwise.
  ℹtem For each session, declare a \textbf{Situation Timer } [4−−8] for the immediate political question.
  ℹtem Play scenes. Any \textbf{1s generate SB}; the GM uses \emph{Social SB options}.
  ℹtem Between sessions, run the \textbf{Intrigue Cycle} to advance timers off−screen and seed the next scandal.
\end{enumerate}

\tableofcontents
≠wpage

\chapter{Political Intrigue – The War of Laws and Lances}
łabel{ch:political−intrigue}

\epigraph{The throne is empty. The cantons sharpen their quills. The condottieri sharpen their swords. And the Everflame burns both.}{— Dusana of the Raven Road, on the Belworth crisis}

This expansion provides procedures for long‑form political conflict in the fractured realms of Vhasia, Viterra, and the contested marches. Here, a duel of advocates can decide what a chevauchée cannot; a printed broadsheet can ruin a bloodline faster than a cannon; and a loan from a coin‑house may buy a crown—or sell it.

\section*{The Three Pressures}

Three forces grind against one another in every court, every hedge, every mercenary camp:

\begin{description}
    ℹtem[The Chivalric Remnant] – Knights who still speak of honour, yet fight for ransom. Jousts are theatre; the tilt is a negotiation. The old codes buckle under the weight of gunpowder and gold.
    ℹtem[The Hedge‑Law Ascendancy] – Magistrates, advocates, and clerks who believe that a well‑written charter is worth a company of lances. Parlements multiply; precedents are printed; the very soil is argued over by the rod.
    ℹtem[The Condotta Engine] – Mercenary captains who sell steel to the highest bidder, and the coin‑houses that bankroll them. Their loyalty is a season long; their ledger is written in blood and interest.
\end{description}

None of these forces can win alone. A knight needs a lawyer to seal the treaty. A magistrate needs a captain to enforce the writ. A condottiere needs a coin‑house to pay the muster. Their friction is your lever.

\section*{What This Expansion Adds}

\begin{itemize}
    ℹtem \textbf{Faction Sheets} with Influence, Stability, and Exposure timers.
    ℹtem \textbf{Intrigue Actions} (Petition, Bribe, Blackmail, Expose, Shield) and a Social SB menu.
    ℹtem The \textbf{Parlement Mini‑Game} – a three‑beat legal duel for courts and councils.
    ℹtem \textbf{The Round of Five Virtues} – a joust scoring system that turns combat into moral theatre.
    ℹtem \textbf{Campaign Dials} (Mandate and Crisis) that track your party’s public legitimacy and the opposition’s growing strength.
    ℹtem \textbf{The Intrigue Cycle} – off‑screen faction moves between sessions.
\end{itemize}

\section*{A Mantle for the Table}

When the game turns from dungeons to duchies, repeat this truth:

\medskip
\noindent\textit{``The lance is judged, but the ledger is weighed. Raise your banner, but keep your seal wet. Winter is coming—and the law does not freeze.''}

\medskip
The following pages give you the tools to play that winter.

\section{Faction Play: Components & Templates}
\subsection{Faction Sheet (Template)}
\begin{verbatim}
[FACTION NAME]
Tier: (street / guild / city / crown / synod)
Aim: (what they want this season)
Stance toward PCs: (Allied / Wary / Hostile)
Levers (3): (money, law, clergy, blades, secrets, ships, votes, rites...)
Leader & Lieutenants: (names + one sentence intent)
Reputation Tags (2): (Merciful, Ruthless, Old−Blood, Reformers, Ascetic...)
Timers:
− Influence [6]: how aligned they are with PC aims; fills to grant access/patronage
− Stability [6−−8]: how coherent the faction remains; fills → schism/coup/purge
− Exposure [6]: public risk & scandal; fills → scandal event triggers
Holdings/Assets: (2−−5 named assets with a sentence of function)
Strings: (1−−3 hooks tying them to locales, laws, rites, or other factions)
\end{verbatim}

\paragraph{When to tick (with examples)}
\begin{itemize}
  ℹtem \textbf{Influence:} Tick when PCs advance or protect the faction's \emph{stated Aim} or humiliate an \emph{active rival} of that Aim. (e.g., secure a river charter the Harbor Guild wants = +1; discredit the Tide−Chart Office opposing them = +1.) Routine gifts without movement on the Aim do not tick.
  ℹtem \textbf{Stability:} Tick when internal strains surface beyond normal bickering: succession wobbles, ideological schism, budget shortfalls, a key lieutenant defects, or magical overreach taxes rites/logistics. (e.g., failed purge = +1; whispered heresy in the Ember−Temple = +1.)
  ℹtem \textbf{Exposure:} Tick on visible missteps: public concessions/humiliations, failed high−profile actions (especially with Social SB spends), leaks of leverage, or collateral harm to bystanders. (e.g., bribery discovered = +2; foiled blackmail exposed = +1.)
\end{itemize}

⫕section*{Strings −−− What They Do}
A \textbf{String} is a concrete tie: ``Controls river tolls,'' ``Holds three Synod votes,'' ``Custodian of the Lantern Rites,'' ``Patron of the Playhouse.''
\begin{itemize}
  ℹtem Acting \emph{with/through} a relevant String: improve \textbf{Position +1} \emph{or} treat \textbf{DV as −−1} (once per scene).
  ℹtem Acting \emph{against} a target's String: \textbf{Position −−1} \emph{or} GM banks +1 SB.
  ℹtem Cutting or seizing a String is a \textbf{4−−6 segment timer}; on fill, move the String to the victor or strike it.
\end{itemize}

\subsection{Relationship Map (Strings & Stakes)}
Sketch the web. For every arrow between two factions, write a \textbf{stake} (e.g., ``river tolls,'' ``chapterhouse votes,'' ``forbidden codex''). Mark two arrows \textbf{volatile}; these will pop first under pressure.

\paragraph{Volatility Triggers}
\begin{itemize}
  ℹtem When either linked faction's \textbf{Exposure} fills, immediately trigger the volatile stake as an on−screen complication or new Situation timer.
  ℹtem When \textbf{Crisis} rises by 1, the GM may trigger one volatile arrow anywhere on the map.
  ℹtem A targeted \textbf{Social SB} spend may flip a volatile stake early (GM choice; announce loudly).
\end{itemize}

\paragraph{Map−to−Action Hooks}
\begin{itemize}
  ℹtem Acting along an arrow where you have \textbf{positive Strings} on either node: \textbf{+1 Position} or \textbf{DV −−1}.
  ℹtem Acting against a node with \textbf{negative entanglement} (Hostile stance or disgraced tag): \textbf{−−1 Position} or GM banks +1 SB.
  ℹtem These are caps; do not chain multiple arrows for stacking.
\end{itemize}

\subsection{Campaign Dials}
\paragraph{Mandate (0−−6)} Table's public legitimacy and door−opening power.
\paragraph{Crisis (0−−6)} Pressure from rivals, costs, and backlash.
Set both to 2 at start unless fiction says otherwise.

⫕section*{Mechanical Hooks}
\textbf{Mandate}
\begin{itemize}
  ℹtem 0−−1: Cold reception. Public venues start \textbf{Desperate} unless you bring Favors/Strings.
  ℹtem 2−−3: Recognized. Once per social scene, treat Position as \textbf{one step higher} when invoking a relevant Favor or String.
  ℹtem 4−−5: Credible power. Once per session, \textbf{cancel} one Social SB spend targeting you in public venues; public institutions set \textbf{DV −−1} versus your petitions.
  ℹtem 6: Standard−bearers. Start public scenes \textbf{Dominant} unless fiction strongly opposes; your endorsements create a 4−timer \emph{Bandwagon} on targets.
\end{itemize}
\textbf{Crisis}
\begin{itemize}
  ℹtem 0−−1: Calm waters.
  ℹtem 2−−3: Tension. At the start of each session, GM banks +1 SB for political scenes.
  ℹtem 4−−5: Boiling. First failed social action each session ticks party \textbf{Exposure +1}.
  ℹtem 6: Breaking point. Start political scenes with a \textbf{Counter−Timer [4]} in play.
\end{itemize}

\section{Scene Procedures for Social Conflict}
\subsection*{Social Conflict at a Glance}
\begin{enumerate}
  ℹtem \textbf{Intent} $→$ \textbf{Position/DV} $→$ \textbf{Roll} $→$ \textbf{Outcome} $→$ \textbf{Ripple}.
  ℹtem Set Dominant/Controlled/Desperate and DV from fiction. \textbf{Any 1s generate SB} in social scenes.
\end{enumerate}

\subsection{Intrigue Actions (examples)}
Each uses the table's attribute/skill combination. Suggested positions/DVs and on−hit effects below.
\begin{itemize}
  ℹtem \textbf{Petition (Presence+Sway):} Present a case to a decision−maker. Dominant with writs/favors; Desperate when hostile. DV 2−−4. Hit: fill Situation; strong hit: tick ally Influence.
  ℹtem \textbf{Bribe (Wits+Streetwise/Lore):} Move resources to open doors. Dominant; worsens if watched/holy ground. DV 2−−3. Hit: create a temporary tag (\emph{Greased Palms}); miss: tick Exposure.
  ℹtem \textbf{Blackmail (Wits+Investigation):} Apply leverage. Dominant $→$ Desperate if unverified. DV 3−−5. Hit: concession now or place a 4−timer \emph{Compromise} usable later for Position Dominant.
  ℹtem \textbf{Broker (Presence+Diplomacy):} Reconcile two parties. Dominant on neutral ground; Dominant with ritual hospitality; Desperate if blades present. DV 3−−4. Hit: reduce both parties' Exposure by 1; fill Influence for the favored side.
  ℹtem \textbf{Expose (Presence+Performance):} Publicly reveal truth/forgery. Dominant with corroboration; Desperate with hearsay. DV 3−−5. Hit: add \emph{Disgraced} tag; fill 2−−3 segments on target Exposure.
  ℹtem \textbf{Shield (Wits+Tactics/Lore):} Pre−empt rival plays via law/protocol/counter−messaging. Dominant; Dominant with counsel/rites. DV 3. Hit: reduce incoming Social SB by 1 this scene; mark ally Stability $−$1 (cohere under pressure).
  ℹtem \textbf{Host Feast/Vigil (Presence+Performance/Rite):} Gather elites; set tone. Dominant if traditions observed; Desperate if taboo. DV 2−−4. Hit: create Audience tags; if triumph, \emph{Mandate +1}.
  ℹtem \textbf{Infiltrate (Body/Wits+Stealth/Deception):} Plant evidence, steal ledgers, ghost a meeting. Desperate when warded/sanctified. DV 3−−5. Hit: gain \emph{Leverage (1)} or fill Situation; on partials, start a \emph{Security} or \emph{Rumor} counter−timer.
\end{itemize}

\subsection{Social SB Options}
When 1s appear, GM spends SB to:
\begin{itemize}
  ℹtem Trigger faux pas; spread rumor; demand concessions; rival interjects; gallery turns; patronage flips.
  ℹtem \textbf{Currency hits:} lose 1 Favor with an implicated institution; a used Leverage becomes public (tick either party's Exposure +1); shift an Audience/Institution tag against the PCs.
\end{itemize}

\subsection{Resolution}
A scene resolves when the \textbf{Situation Timer } fills or fiction dictates. Then ripple:
\begin{itemize}
  ℹtem Advance/reduce involved \textbf{Faction Timers} (Influence/Stability/Exposure).
  ℹtem Adjust \textbf{Mandate/Crisis} (usually $±$1 each based on victory and visible costs).
  ℹtem Record lasting \textbf{Audience/Institution Tags}.
\end{itemize}

\subsection{Situation vs. Faction Timers −−− Default Ripple}
\begin{itemize}
  ℹtem \textbf{Success (PC−favoring):} Allied \emph{Influence +1} (or +2 if it advances their Aim); if win was clean, allied \emph{Exposure −−1}; hostile \emph{Exposure +1}.
  ℹtem \textbf{Failure (PC−hurting):} Opposing \emph{Influence +1} (or +2 if on−Aim); allied \emph{Exposure +1}. If failure was internal, also tick that faction's \emph{Stability +1}.
  ℹtem \textbf{Compromise:} Each side \emph{Influence +1} or both \emph{Exposure +1}, guided by fiction.
\end{itemize}

\section{The Intrigue Cycle (Between Sessions)}
Use once per session (or per in−fiction week/month on jumps).
\subsection*{Step A −−− Update World State}
For each active faction, advance \emph{Stability} or \emph{Exposure} by 1 if they suffered public pressure or internal strain. Tick unresolved thread timers by 1 or start a 4−−6 timer.

\subsection*{Step B −−− Draw/Prompt for Events}
\textbf{Method:} Draw blindly from a standard deck or deliberately pick suits matching themes. Use \textbf{Suit} for type (Hearts=court/crowd; Spades=violence; Clubs=resources/strikes; Diamonds=environment/omens) and \textbf{Rank} for severity (low=+1, mid=+2, face/ace=+3 segments).\\
\textbf{Apply:} Place each card on a faction/timer and immediately tick by severity; write one sentence of fallout that must surface next session.\\
\textbf{Examples:} Hearts 7 $→$ public oath/romance, tick \emph{Exposure +2}. Clubs 10 $→$ trade shock/strike, start a 6−timer \emph{Shortages}.

\subsection*{Step C −−− Off−Screen Actions}
Each PC directs one off−screen move using a follower/asset/network. Resolve with one roll at \emph{Significant Time} scale; on 1s, bank SB and/or tick Exposure.

\subsection*{Step D −−− Recompute Dials}
If wins were visible and valued, \textbf{Mandate +1}. If costs, scandals, or riots mounted, \textbf{Crisis +1}. Both may move.

\subsection*{Step E −−− Prep Next Situation}
Name the next immediate political question. Put a \textbf{4−−8} timer on it. Grab 2−−3 complications tied to factions and prior tags.

\section{Currencies: Favor, Leverage, Exposure, Reputation}
\subsection{Favor (Table Currency)}
Earn a Favor when you deliver visible benefit or save a leader's face. Spend a Favor to:
\begin{itemize}
  ℹtem Improve Position by one step \emph{against that person/institution}.
  ℹtem Cancel one Social SB spend aimed at you in their venue.
  ℹtem Access a restricted office/archive/sanctum once.
\end{itemize}
Favors are narrow; specify scope (``Harbormaster's Office,'' not ``All of Kahfagia''). Three unspent favors from the same seat convert to \textbf{Standing (Tag)} with them.

\subsection{Leverage (Portable)}
\begin{itemize}
  ℹtem One \textbf{Leverage (1)} forces a concession (1 segment on a Situation Timer ) \emph{or} holds Position at Dominant when it would drop to Desperate.
  ℹtem \textbf{Decay:} Only \emph{public} uses count. After two public uses of the same leverage, convert it to \emph{Exposure +1} on you or the target (your choice). Private uses do not decay it, but three private uses make it stale; refresh with new proof/witness/circumstance before it functions again.
\end{itemize}

\subsection{Exposure (You & Them)}
Track \textbf{Exposure [6]} for the party (or per face) and on major factions. When a track fills, trigger a \textbf{Scandal Event}: public censure, hostile audit, inquisitorial summons, riot, duel demand, etc. Clear 2 on a successful \emph{Shield} scene or by accepting a humiliating concession on−screen.

\subsection{Reputation Tiers}
Check at arc end (or every 2−−3 sessions). If two are true, advance Reputation Tier:
\begin{itemize}
  ℹtem Convert 3 favors from one institution into \textbf{Standing (Tag)} \emph{or} hold Standing with two institutions.
  ℹtem End with \textbf{Mandate $≥$ 4} \emph{or} resolve a city−scale Situation Timer [8] in your favor.
  ℹtem Your deeds create/resolve a \emph{regional law/rite} or seat/remove a leader.
\end{itemize}
\textbf{Gatekeeping Effects:} Higher tiers grant default access to higher−tier factions/venues; DVs for petitions to lower−tier offices drop by 1; timers scale up (use 6−−8 segments at higher tiers; 4−−6 at low tier). When your tier outstrips a scene's stakes, the GM may simply hand you the room unless opposing Strings/Leverage are present.

\section{Followers & Assets in Politics}
\subsection*{On−Scene}
Assign a follower to Assist or to take an independent scene action with its own risk. Harm to followers often lands as \emph{Stability} or \emph{Exposure} ticks on their home faction.

\subsection*{Off−Screen Orders (Intrigue Cycle)}
Each PC may direct one follower/asset on a \emph{Significant−Time} operation (surveil, seed rumors, counter−bribe, audit a ledger, escort witnesses). This can occur mid−adventure as a cutaway or between sessions. Resolve with one roll using the follower's capability as positioning; on a hit, apply 1−−2 segments to the intended timer; on 1s, bank SB and mark \emph{Exposure +1}.

\subsection*{Upkeep & Entanglements}
Track names, not ledgers. When the GM spends 2+ SB in a political arc, an ally may face the consequence instead (compromised informant, bribed clerk, stolen seal). Convert entanglements into new \emph{threads} and \emph{timers} rather than binary losses.

\section{Generators (Scandal Decks & Prompts)}
\subsection*{Suits as Political Color}
Hearts: court favor, marriages, oaths, crowd mood.\\
Spades: violence, duels, arrests, street pressure.\\
Clubs: resources, trade shocks, supply, strikes.\\
Diamonds: environment, festivals, omens, disasters.

\subsection*{Ranks as Severity}
Low (2−−5): gossip/minor obstruction ( +1 segment).\\
Mid (6−−9): scandals, votes, riots, inquisitors ( +2 segments).\\
Face/Ace: institutions shift; laws change; thrones wobble ( +3 segments; start/finish a major timer).

\subsection*{Draw Procedures}
Quick Hook (2): what's the scandal, who is entangled?\\
Full Seed (4): place each onto a faction/timer; one−sentence fallout per card.\\
Act Builder: Establishment $→$ Escalation $→$ Resolution; note which timers must move.

\subsection*{Worked Examples}
\begin{itemize}
  ℹtem \textbf{Hearts 7:} A scion's oath of love becomes public; tick \emph{Exposure +2} on their house; start a 4−timer \emph{Secret Betrothal}.
  ℹtem \textbf{Spades King:} A duelist challenges the Tide−Marshall; start a 6−timer \emph{Public Blood}; if it fills, \emph{Mandate −−1} for the city guard.
  ℹtem \textbf{Clubs 6:} Dockworkers strike; tick Harbor Guild \emph{Stability +1}; create tag \emph{Pickets at Dawn}.
  ℹtem \textbf{Diamonds 9:} Omens sour during a procession; tick \emph{Exposure +2} on the Ember−Temple unless rites appease it.
\end{itemize}

\section{Long−Term Campaign Management}
Running a multi−arc intrigue campaign requires pacing and player investment. Use these additional tools.

\subsection*{Seasonal Arcs}
Each arc (4−−6 sessions) should focus on a single \textbf{Seasonal Aim} for the party (e.g., ``Seat an allied candidate'', ``Expose a treasonous conspiracy'', ``Negotiate a trade treaty''). At the end of the arc, the party may:
\begin{itemize}
  ℹtem Permanently increase \textbf{Mandate} by 1 (if the aim succeeded and was visible).
  ℹtem Gain a \textbf{Legacy String} (a permanent tie usable once per arc).
  ℹtem Convert a major rival faction's \textbf{Influence} timer into a \textbf{Debt} on them.
\end{itemize}

\subsection*{Faction Rivalry Timers}
In addition to Influence/Stability/Exposure, you may track a \textbf{Rivalry Timer [6]} between two factions. Each time one faction publicly humiliates the other, tick the timer. When full, a decisive confrontation occurs (duel, vote, raid, trial). The PCs can influence the outcome by spending Favors or Leverage.

\subsection*{Managing Downtime}
Between arcs, run a \textbf{Downtime Phase} where players narrate their characters' activities (3−−5 scenes). Each activity:
\begin{itemize}
  ℹtem Restores a resource (Favor, Leverage, or clears Exposure).
  ℹtem Advances a personal project (gaining a new String or improving an Asset).
  ℹtem May trigger a \textbf{Personal Complication} (GM banks 2 SB for the next arc).
\end{itemize}

\subsection*{Delegation While Away – Maintaining the Household in Absentia}
łabel{subsec:delegation−away}

When the party is absent for a full season (spring, summer, autumn, or winter) on an adventure, diplomatic mission, or military campaign, the household does not simply pause. All timers would normally tick –1 at season end, but the party may avoid a death spiral by delegating authority to a trusted follower.

⫕section*{Core Principle}
A \textbf{Cap 4 or Cap 5 follower} (see \textit{Player's Guide}, §12.4) who has been delegated a minimum number of \textbf{maintained assets} (see \textit{Player's Guide}, §12.14.1.1) may act as Steward. They prevent one or two household timers from automatically ticking and may attempt to improve those timers using their own capabilities.

⫕section*{Asset Requirements}
The follower must have at least the following number of assets \textbf{delegated to them and in Maintained condition}:

\begin{center}
\begin{tabular}{|c|c|c|}
ℎline
\textbf{Follower Cap} & \textbf{Minimum Maintained Assets Delegated} & \textbf{Timers They Can Manage} \\
ℎline
4 & 1 & 1 \\
5 & 2 & 2 \\
ℎline
\end{tabular}
\end{center}

\noindent
A Cap 4 follower with two assets is still limited to one timer – the extra asset does not grant a second. Only Cap 5 followers can manage two timers simultaneously (e.g., Loyalty and Security).

⫕section*{Procedure}
Before departure, the party:
\begin{enumerate}
    ℹtem Names a qualified follower as \textbf{Steward}.
    ℹtem Designates which household timer(s) the Steward will manage (up to the limit above).
    ℹtem The Steward remains at the keep for the entire season.
\end{enumerate}

At the end of the season:
\begin{enumerate}
    ℹtem All \textbf{non‑delegated} timers tick –1 (wear of time).
    ℹtem For each delegated timer, the Steward attempts to \textit{improve} it using the normal \textbf{Household Action} rules (see §\ref{sec:household−actions}). The Steward rolls a number of d10s equal to their \textbf{Cap}. They do not add attributes or skills unless they have a relevant talent (e.g., \textit{Efficient Steward}). The GM sets Difficulty Value (DV) and Position based on the Steward’s competence, the season’s fiction, and the state of the holding.
    ℹtem Apply the standard outcome matrix:
    \begin{itemize}
        ℹtem \textbf{Clean Success:} Timer advances by 1.
        ℹtem \textbf{Success with SB:} Timer advances by 1, GM gains 1 Story Beat (complication).
        ℹtem \textbf{Partial:} Timer unchanged, but the Steward learns a secret (gain a String) or the GM gains 1 SB.
        ℹtem \textbf{Miss:} Timer ticks backward by 1 (or a minor crisis occurs).
    \end{itemize}
    ℹtem If a delegated timer reaches 0 despite the Steward’s efforts, a \textbf{Crisis Event} (see §\ref{sec:crisis−events}) occurs. The crisis must be resolved by the party upon their return; the Steward cannot resolve it alone.
\end{enumerate}

⫕section*{Limitations and Strategic Choices}
\begin{itemize}
    ℹtem Only one follower may act as Steward per season. (If multiple qualified followers remain, they may assist, but only the designated Steward rolls.)
    ℹtem Cap 3 or lower followers cannot delegate – they lack the authority or infrastructure to manage a noble household alone.
    ℹtem If two or more player characters remain at the keep, they may each take a household action normally; delegation is only needed when \textbf{all} player characters are absent.
    ℹtem A Steward who is also the player’s heir or successor (see \textit{Player's Guide}, Legacy Engine) is an excellent narrative choice – delegation becomes part of their training.
\end{itemize}

⫕section*{Optional Rule – Delegation with Risk}
For tables that enjoy deeper strategic planning, players may give the Steward \textbf{specific orders} before departure. The player rolls \textit{Presence + Command} (or a relevant skill) against a DV set by the complexity of the order (usually 3–4). On a success, the Steward gains +1 die on their season‑end roll for each delegated timer. On a failure, the Steward is confused or resentful, and each delegated timer ticks –1 regardless of the Steward’s roll (in addition to any other outcome).

This optional rule adds a layer of pre‑planning but also extra dice rolls; use it only if your table enjoys such strategic depth.

⫕section*{Integration with Existing Rules}
\begin{itemize}
    ℹtem This rule assumes the four standard household timers (\textit{Loyalty, Prosperity, Security, Hospitality}). If additional timers are used (e.g., \textit{Armor Closes} for Vhasia, \textit{Legal Exposure} for Viterra), the party may delegate \textbf{one additional timer per extra Cap 5 follower} (maximum two delegated timers per season, requiring two separate Stewards).
    ℹtem See \textit{Player's Guide} §12.4 for follower Cap costs (Cap 4 costs 16 XP; Cap 5 costs 25 XP). The asset requirements ensure that delegation is a reward for investing in both the follower and their infrastructure.
\end{itemize}

\chapter{Court & Keep – Managing a Noble Household}
łabel{ch:courtkeep}

\epigraph{A keep is not stone. It is a ledger of loyalties, a harvest of favours, a winter of patience.}{— Duchess Gwen of the Thornwood, on her sixtieth year of rule}

This chapter provides procedures for players who hold a noble title, command a fortress, or seek to build a political dynasty. It is designed for campaigns that span seasons or years, where influence, heirs, and household management matter as much as swords and spells.

\section{The Noble Household – Core Components}

A noble household consists of four interconnected tracks. Each is a timer (size 6−8) that the party must maintain. Neglect one, and the others may suffer.

\begin{description}
    ℹtem[Loyalty [6]] – The devotion of your household staff, lesser nobles, and common tenants. High loyalty means servants warn you of plots; low loyalty means whispers in dark corners.
    ℹtem[Prosperity [6]] – Harvests, trade, coin reserves, and infrastructure. High prosperity funds improvements and bribes; low prosperity leads to debt, famine, and desertion.
    ℹtem[Security [6]] – Walls, patrols, wards, and spy networks. High security deters raids and assassins; low security invites them.
    ℹtem[Hospitality [8]] – Your reputation as a host, the splendour of your feasts, the loyalty of visiting knights and envoys. High hospitality attracts allies; low hospitality offends potential patrons.
\end{description}

\subsection*{Seasonal Upkeep}
At the end of each season (spring, summer, autumn, winter), each timer ticks automatically by 1 (the wear of time). The party may take actions to reduce or reverse these ticks (see below). If any timer reaches 0, a \textbf{Crisis Event} occurs – see the Crisis table.

\section{Running the Household – Actions Per Season}

During each season, the party may assign themselves or trusted followers to one or more of the following actions. Each action requires a roll (or, if using a follower, a use of their capability). A character can only perform one action per season unless they have a relevant Asset or Talent (e.g., \textit{Efficient Steward}).

\subsection*{Improve Loyalty (Presence + Command)}
Roll to raise the loyalty of a specific group (servants, garrison, tenant farmers). DV 3 (peaceful region) to 5 (discontent after a harsh winter).

\begin{itemize}
    ℹtem \textbf{Success:} Loyalty timer advances by 1.
    ℹtem \textbf{Success with SB:} Advances by 1, but GM gains 1 SB (rival spreads a rumour).
    ℹtem \textbf{Partial:} No change, but you learn a secret (String).
    ℹtem \textbf{Miss:} Loyalty ticks backward by 1 (a trusted servant resigns).
\end{itemize}

\subsection*{Improve Prosperity (Wits + Craft, Commerce, or Lore)}
Invest in trade, farming, or infrastructure. DV 3 (normal year) to 5 (drought or siege).

\begin{itemize}
    ℹtem \textbf{Success:} Prosperity advances by 1. You may also gain a minor Asset (e.g., new mill, trade route).
    ℹtem \textbf{Partial:} Prosperity advances by 1, but you incur a Debt Timer [4] (loan from a merchant or patron).
    ℹtem \textbf{Miss:} Prosperity ticks backward (blight, embezzlement).
\end{itemize}

\subsection*{Improve Security (Wits + Tactics, Stealth, or Arcana)}
Patrols, espionage, or warding. DV 3 (peace) to 5 (border war).

\begin{itemize}
    ℹtem \textbf{Success:} Security advances by 1. Reveal one hidden threat (an assassin, a spy, a cursed object).
    ℹtem \textbf{Partial:} Security advances by 1, but the GM gains 1 SB (the enemy learns your methods).
    ℹtem \textbf{Miss:} Security ticks backward; a minor breach occurs (stolen letter, dead sentry).
\end{itemize}

\subsection*{Improve Hospitality (Presence + Performance or Sway)}
Host a feast, a hunt, a tourney, or a diplomatic reception. DV 3 (modest) to 5 (lavish, with many guests).

\begin{itemize}
    ℹtem \textbf{Success:} Hospitality advances by 1. Gain a temporary String with one guest faction.
    ℹtem ℹtem \textbf{Partial:} Hospitality advances by 1, but the feast incurs a minor scandal (drunken brawl, insulted herald). Tick Exposure +1.
    ℹtem \textbf{Miss:} Hospitality ticks backward; the feast is a disaster (poisoned wine, open duel). Crisis may trigger.
\end{itemize}

\section{Reputation Across Seasons – The Renown Dial}

In addition to the four timers, the party has a \textbf{Renown Dial} ranging from –3 (infamous) to +3 (exalted). Renown affects all social rolls with nobles, traders, and foreign courts.

\begin{itemize}
    ℹtem \textbf{−3 to −1:} Disgraced. You are received Desperate in any formal setting. Trade goods cost double. Followers may desert.
    ℹtem \textbf{0:} Neutral. Standard Position. No modifier.
    ℹtem \textbf{+1 to +3:} Respected. You start social scenes Controlled or Dominant. You may demand audience with higher nobles. Merchants offer credit (one free Debt Timer of 4 segments per season).
\end{itemize}

\subsection*{Changing Renown}
Renown shifts by ±1 at the end of each season based on:

\begin{itemize}
    ℹtem \textbf{Major victory in war or law:} +1
    ℹtem \textbf{Hosting a legendary feast (Hospitality filled to 8):} +1
    ℹtem \textbf{Public scandal or failed oath:} –1
    ℹtem \textbf{Losing a keep or being defeated in a duel:} –1
    ℹtem \textbf{Marrying into a powerful house:} +1
\end{itemize}

The GM may also award or deduct Renown based on story events.

\section{Marriage Alliances – Skill Challenge}

Marriage is the sharpest tool of diplomacy. To arrange a marriage between your house and another, run a skill challenge. The challenge requires \textbf{3 successes before 2 failures}. Each success represents a key hurdle overcome; each failure introduces a complication that may raise the price.

\subsection*{Hurdles (examples)}
\begin{itemize}
    ℹtem \textbf{Proposal (Presence + Command):} Formally ask the other house for the match. DV 3 if they are eager; DV 5 if they are reluctant.
    ℹtem \textbf{Dowry or Bride−Price (Wits + Commerce or Lore):} Negotiate the exchange of lands, gold, or favours. DV 4.
    ℹtem \textbf{Religious Approval (Spirit + Lore):} Secure a priest's blessing, or navigate differing faiths. DV 3−4.
    ℹtem \textbf{Convince the Betrothed (Presence + Sway):} The prince or princess may have their own opinions. DV 4.
    ℹtem \textbf{Feast (Presence + Performance):} Announce the betrothal before witnesses. DV 3.
\end{itemize}

\subsection*{Outcomes}
\begin{itemize}
    ℹtem \textbf{3 Successes:} The marriage is agreed. Gain a permanent Alliance String with the other house. Increase Renown by +1. The couple produces an heir (a new NPC).
    ℹtem \textbf{2 Successes, 2 Failures (partial):} The marriage goes forward, but at a cost. Increase Renown by +0. Gain a String, but also a Debt Timer [6] (the other house expects a future favour).
    ℹtem \textbf{Failure (2 failures before 3 successes):} The marriage fails. Renown –1. The other house may become hostile (new Rival Timer [8]).
\end{itemize}

\subsection*{Example: The Betrothal of Lord Reynard and Lady Celandine}
The party acts as Lord Reynard's advocates. They roll:
\begin{enumerate}
    ℹtem \textbf{Proposal (Presence+Command, DV 4):} Success – 1 success.
    ℹtem \textbf{Dowry (Wits+Commerce, DV 5):} Partial – no success, but no failure. GM gains 1 SB.
    ℹtem \textbf{Convince the Betrothed (Presence+Sway, DV 4):} Failure – 1 failure.
    ℹtem \textbf{Religious Approval (Spirit+Lore, DV 3):} Success – 2 successes, 1 failure.
    ℹtem \textbf{Feast (Presence+Performance, DV 3):} Success – 3 successes, 1 failure. \textbf{Partial success.} The marriage is agreed, but the dowry negotiations were strained; the party incurs a Debt Timer [6] (a future political favour).
\end{enumerate}

\section{Crisis Events (When a Timer Reaches 0)}

When any household timer reaches 0, roll or choose a Crisis Event from the table below. The crisis is immediate and requires a scene to resolve. After resolution, the timer resets to 2 (or 3 for Hospitality), and the party may begin rebuilding.

\begin{longtable}{|p{0.2\textwidth}|p{0.6\textwidth}|}
ℎline
\textbf{Clock} & \textbf{Crisis Event} \\
ℎline
Loyalty & Servants betray a secret to a rival; a faction of tenants threatens to leave. To restore, the party must perform a public act of generosity or punish a known abuser. \\
ℎline
Prosperity & Famine or bankruptcy. Food stores empty; creditors demand immediate payment. The party must find a new source of grain (quest) or sell an Asset. \\
ℎline
Security & A raid, an assassination attempt, or a magical breach. The party must defend the keep (combat skill challenge) or track down the infiltrator. \\
ℎline
Hospitality & A guest is insulted, injured, or murdered under your roof. Your reputation crumbles. The party must offer weregild, host a second feast, or duel the offended party. \\
ℎline
\end{longtable}

\section{Optional Rule – Heirs and Succession}

If the party wishes to establish a dynasty, track the primary heir as an NPC. At the end of each year (four seasons), roll a die:

\begin{itemize}
    ℹtem \textbf{1−2:} The heir falls ill or suffers an accident (timer [4] to recover).
    ℹtem \textbf{3−5:} No event.
    ℹtem \textbf{6:} The heir comes of age. The player may transfer their character to the heir (retaining some skills and Strings) or keep the heir as a powerful follower (Cap 4−5).
\end{itemize}

The heir inherits the household's current timers (adjusted for the new generation) and gains one bonus Talent based on their upbringing (e.g., \textit{Disciplined Body} for martial education, \textit{Courtly Graces} for diplomatic).

\section{GM Guidance}
The \textbf{Court & Keep} framework is designed for campaigns that value long−term consequences. Use it when players acquire a fortress, a title, or a permanent base. Emphasise that each season brings new opportunities and threats. Let the players choose which actions to prioritise – they cannot fix everything at once.

A simple mantra for the table: \textit{``Stone and title are debts. Pay them on time, or the winter will collect.''}
\clearpage

% ===================================================================
%   BLOOD ON THE BELWORTH ADVENTURE
% ===================================================================

\begin{center}
    {\Huge \textbf{Blood on the Belworth}} \\
    \vspace{0.5em}
    {Łarge A Political Intrigue adventure for Fate's Edge using the v0.3 framework}
\end{center}

≤ttrine{D}{uke} Braedon Fenwood, the last of the old Valvano line, has just died. His sons −−− the cautious Constano (heir) and the martial Seniro −−− now struggle to maintain House Fenwood's power on both sides of the Belworth River. The succession crisis of Viterra festers: Queen Veralyn Demoneye's throne is contested, and the Vhasian Regency sees opportunity. In the midst of this, a ferry is found adrift with three stabbed corpses and a Fenwood quartermaster's armband nailed to the mast. The river turns to rumor, and the old duke's final words echo: \textit{"Count exits before enemies. The road reschedules."}

\section*{Timeline Note (GM Eyes Only)}
The Fourth Sack of Ecktoria (c. 620 AR) occurred seventy years after the Third Sack and the Valvano exile. The events of \textit{The Exile Road} and \textit{The Nameless} take place around that era. Braedon Fenwood lives centuries later, after the Tarling kings have fallen, the Fenwoods have risen to ducal status, and the Everbloods have taken the crown. The debt foreshadowed in \textit{The Exile Road} was paid long ago. This adventure deals with the \textit{consequences} of those ancient choices: the Black Scrip indemnities, the ghost of the Thornhill child, and the Aeler keystone rights.

\section*{Using This Adventure}
This module is built on the Political Intrigue Campaign Framework (v0.3). Run each major question with a Situation Timer (4−−8) and at least one Counter−Timer (4−−6). Use Strings to tune Position/DV and as prizes/failure states. Apply Mandate/Crisis dials and the Situation→Faction Ripple rules after each scene.

\textbf{Intended length:} 2−−4 sessions. \textbf{Tone:} border−court drama, false flags, rites of hospitality, hedge−law wrangling. \textbf{Prerequisite:} The adventure assumes the events of \textit{The Old Duke's Wars} (vignettes of Braedon's life) have recently occurred. The party may have played those vignettes, or the GM can narrate the key beats.

\section*{Seats, Locales & Assets}
\begin{itemize}
  ℹtem \textbf{Tarlington (Viterra):} Ducal seat; marble halls, Hedge−Law Moot Court, lantern beacons along the river stairs. Old Braedon's solar now holds his Canray board and a sealed letter.
  ℹtem \textbf{Belworth River Corridor:} Ferry towns, reedmarshes, customs towers, lantern−law posts, ruined Tarling fortlets.
  ℹtem \textbf{Lence (Vhasia):} Half−rebuilt after the sack; Restorationist salons; a palace−in−exile flavor.
  ℹtem \textbf{The Red Span:} A half−finished Aeler bridge across the Belworth; a scandal of budgets and graft. The Aeler True Masons have a dormant keystone clause that could swing control.
  ℹtem \textbf{Lantern−Rites House of Oria:} Neutral ground by custom; oaths, feasts, and river burials are sanctified here. Sister Oria was a confidante of old Braedon.
\end{itemize}

\subsection*{Named Assets (examples):}
\begin{itemize}
  ℹtem Ferry Passbooks (who may legally cross at dusk)
  ℹtem Tarler's Seal (ducal customs stamp)
  ℹtem Ledger of Tithes (grain barges)
  ℹtem Lantern−Rite Basin (ritual hospitality)
  ℹtem The Black Scrip (indemnities owed by Lence after the Sack)
  ℹtem Red Span Budget Roll (inflated contracts)
  ℹtem \textbf{Braedon's Canray Board} (a slate with an unfinished game; the positions reveal hidden alliances)
\end{itemize}

\section*{Factions (Sheets)}

\subsection*{House Fenwood (Viterra)}
\begin{longtable}{@{}lll@{}}
Tier: crown & Aim: Preserve cross−river power and ensure a smooth succession for Veralyn & Stance: Needs the PCs \\
Levers: blades, law, tolls, Canray wisdom & Leader: Duke Braedon (deceased); Constano (heir), Seniro (marshal) & \\
Reputation: Old−Blood, Stern & & \\
Timers: Influence [6], Stability [8], Exposure [6] & & \\
Strings: Custodian of the Lantern Beacons; Controls Red Span work crews; Holds oaths of three ferry captains; \textit{Braedon's unfinished game −− a metaphor for the succession} & & \\
\end{longtable}

\subsection*{Vhasian Regency Council}
\begin{longtable}{@{}lll@{}}
Tier: crown & Aim: Cement a successor loyal to Lence and challenge Fenwood indemnities & Stance: Wary \\
Levers: law, clergy, spies, ancestral grievances & Leaders: Lady−Envoy Mael Vesht; Arch−Rector Irsen & \\
Reputation: Restorers, Grievant & & \\
Timers: Influence [6], Stability [6], Exposure [6] & & \\
Strings: Custodian of the Black Scrip; Holds exile petitions; Sanctions river pilgrimages & & \\
\end{longtable}

\subsection*{Lantern−Rites Custody (Neutral Clergy)}
\begin{longtable}{@{}lll@{}}
Tier: synod & Aim: Keep rites sovereign and bridge peace & Stance: Allied if respected \\
Levers: rites, hospitality, crowds & Leader: Sister Oria & \\
Reputation: Merciful, Unbending in ceremony & & \\
Timers: Influence [6], Stability [6], Exposure [6] & & \\
Strings: Neutral Sanctum of Hospitality; Funerary authority on the Belworth; Oath−tablet registry & & \\
\end{longtable}

\subsection*{River Syndicate (Pilots & Ferrymen)}
\begin{longtable}{@{}lll@{}}
Tier: guild & Aim: Profit from stability and exemptions & Stance: Opportunists \\
Levers: ships, tolls, rumor, fog−sense & Leader: Captain Pell Rivermark & \\
Reputation: Practical, Superstitious & & \\
Timers: Influence [6], Stability [6], Exposure [6] & & \\
Strings: Ferry Passbooks; Night−lantern routes; Smuggler hush−oaths & & \\
\end{longtable}

\subsection*{Red Span Contractors (Graft Bloc)}
\begin{longtable}{@{}lll@{}}
Tier: merchant & Aim: Keep funds flowing −− finish or delay as profits dictate & Stance: Smiling knives \\
Levers: coin, labor, saboteurs, counterfeit ledger pages & Leader: Master Mason Ghiro & \\
Reputation: Gilded, Cut−Rate & & \\
Timers: Influence [6], Stability [6], Exposure [6] & & \\
Strings: Hold on budget rolls; Union stevedores; Secret pact with Seniro's quartermaster & & \\
\end{longtable}

\subsection*{Lence Restorationists (Radicals)}
\begin{longtable}{@{}lll@{}}
Tier: street→guild & Aim: Public trial of Fenwood war‑crimes; purge collaborators & Stance: Hostile to waffling \\
Levers: crowds, broadsheets, knives, martyrs & Leader: Jaska the Poet & \\
Reputation: Zealous, Wounded & & \\
Timers: Influence [6], Stability [6], Exposure [6] & & \\
Strings: Printing house in the Reeds; Honor cult of the fallen; List of collaborators (real and imagined) & & \\
\end{longtable}

\section*{Relationship Map (Strings & Stakes)}
\begin{itemize}
  ℹtem Fenwood ↔ Regency: \textit{Black Scrip indemnities}; \textbf{volatile}.
  ℹtem Fenwood ↔ River Syndicate: \textit{Ferry exemptions}; \textbf{volatile}.
  ℹtem Regency ↔ Restorationists: \textit{Street muscle vs legitimacy}.
  ℹtem Fenwood ↔ Red Span Contractors: \textit{Overruns & kickbacks}; \textbf{volatile}.
  ℹtem Lantern−Rites ↔ Everyone: \textit{Hospitality oaths}.
  ℹtem Syndicate ↔ Contractors: \textit{Work crews & permits}.
\end{itemize}
\textbf{Volatility Triggers:} Exposure fill on either node; Crisis +1; or a targeted Social SB spend. \\
\textbf{Map Hooks:} Acting with a positive String on a node → Position +1 or DV −−1 (once/scene). Acting against a hostile node → Position −−1 or GM banks +1 SB.

\section*{Campaign Dials (Start)}
\textbf{Mandate:} 2 (recognized but brittle) \\
\textbf{Crisis:} 3 (tension after Brightstride's death and the contested succession)

Move them honestly: public wins/losses change access and pressure as per framework rules.

\section*{Inciting Incident −− "Blood on the Ferry"}
At dusk, the ferry \textit{Belworth Grace} drifts to shore with its lanterns guttering. A Vhasian courier and two clerks lie stabbed; a Fenwood quartermaster's armband is nailed to the mast. Survivors whisper conflicting tales. The river turns to rumor overnight.

In Braedon's solar, his Canray board has been moved −− the Blue piece is now isolated. Those who see it read a warning: \textit{the king is vulnerable}.

\timer{Immediate Situation Timer }{6}: \textbf{Contain the Spiral} −− keep the border calm for 24 hours. \\
\timer{Counter−Clock}{4}: \textbf{Hotheads March} −− Seniro's men assemble to cross at dawn. \\
\timer{Counter−Clock}{4}: \textbf{Justice Now!} −− Restorationists demand arrests at the Lantern−Rites House.

If the PCs do nothing: \textbf{Crisis +1}, Fenwood Exposure +1, Regency Influence +1, and the volatile stake \textit{Black Scrip} triggers (audit squads arrive).

\section*{Act I −− Holding the River}
\textbf{Goal:} Prevent escalation; establish an inquest under Lantern−Rites.

\subsection*{Core Scenes (pick 2−−3)}

\paragraph{Petition at Lantern−Rites}
Claim neutral inquest rights before Sister Oria.  
\position{Dominant with offerings (a candle, a tithe, a memory of Braedon); Desperate if arriving with blades.} \dv{3−−4}  
\textbf{On hit:} Situation +1−−2; Lantern−Rites Influence +1.  
\textbf{On miss:} Exposure +1 (you profane rites).

\paragraph{Shield the Captains}
Keep Seniro's hotheads from marching. Use Hedge‑Law precedents or Sister Oria's writ.  
\position{Dominant with Hedge‑Law dicta; Dominant if you secure Oria's seal.} \dv{3}  
\textbf{On hit:} Counter‑Timer \textit{Hotheads March} −−1; Seniro grudgingly respects you.

\paragraph{Broker a Truce}
Arrange a dusk meeting between Regency envoys and Fenwood stewards at Lantern−Rites.  
\position{Dominant if hosting on sacred ground; Desperate if in the open.} \dv{3−−4}  
\textbf{On hit:} Both sides" Exposure −−1; start a concession timer \textit{Shared Inquest} [4].

\paragraph{Infiltrate the Red Span}
Sneak onto the unfinished bridge to see who planted the quartermaster's armband.  
\position{Desperate (saboteurs may still be present).} \dv{4−−5}  
\textbf{On hit:} Gain Leverage (1) −− \textit{"Seniro's quartermaster took bribes from Ghiro"} −− or fill \textit{Contain the Spiral} by 1.

\subsection*{Vignette Echo (from \textit{The Old Duke's Wars})}
If the party spared the Thornhill child in the vignettes, his ghost appears during infiltration, whispering warnings. This grants +1 die to the roll but adds a \textbf{Haunted} condition (Dread +1).

\textbf{End of Act I:} If \textit{Contain the Spiral} fills, Mandate +1. If any counter‑timer fills first, Crisis +1 and trigger a volatile arrow (GM's choice: likely \textit{Fenwood ⇄ Contractors} or \textit{Fenwood ⇄ Regency}).

\section*{Act II −− Ledger and Blood}
\textbf{Goal:} Identify the hand behind the ferry murders.

\subsection*{Clue Threads (any 3 lead to the culprit)}
\begin{enumerate}
  ℹtem \textbf{Tithe Ledger} −− Missing customs stamps the night of the murder (points to Syndicate colluding with Contractors).
  ℹtem \textbf{Printing House in the Reeds} −− Broadsheets were printed hours after the killings (Restorationists prepped the narrative).
  ℹtem \textbf{Black Scrip Clerk} −− The dead courier carried documents for a renegotiation of the indemnity (someone wanted that negotiation to fail).
  ℹtem \textbf{Quartermaster's Armband Fibers} −− Match a Red Span uniform batch, not Fenwood's armory (Contractors frame‑job).
\end{enumerate}

\subsection*{Likely Culprits (choose one or let players expose one)}
\begin{itemize}
  ℹtem \textbf{Ghiro's Contractors} −− Staged a false flag to prolong funding and tie Seniro closer.
  ℹtem \textbf{Restorationists} −− Struck to collapse any indemnity deal, forcing a hard break.
  ℹtem \textbf{Syndicate Hardliners} −− Provoked crisis to win permanent night‑route exemptions.
\end{itemize}
\textit{Tip: leave two culprits plausible; the third is a red herring that still moves timers.}

\subsection*{Core Scenes (pick 3−−4)}

\paragraph{Expose Budget Graft at the Moot Court}
\position{Dominant with corroboration (ledger + witness).} \dv{4−−5}  
\textbf{On hit:} Target faction gains \textit{Disgraced} tag; Exposure +2; Mandate +1 if public.

\paragraph{Blackmail a Syndicate Captain}
\position{Dominant → Desperate if leverage unverified.} \dv{3−−5}  
\textbf{On hit:} Place \textit{Compromise} [4] on the captain (later shift Position to Dominant once).

\paragraph{Host a Feast or Vigil at Lantern−Rites}
Cool the crowd.  
\position{Dominant if rites observed.} \dv{2−−3}  
\textbf{On hit:} Create Audience Tags (\textit{Warmed}, \textit{Sobering}); Mandate +1 on triumph.

\paragraph{Petition the Regency to Stay Arrests}
\position{Dominant (you have neutral ground).} \dv{3}  
\textbf{On hit:} Tick \textit{Shared Inquest} [4] twice; Regency Influence +1 (they respect your diligence).

\textbf{End of Act II:} If culprit identified on‑screen and public mood calms, reduce \textbf{Crisis −−1} (min 0). On messy wins, tick Fenwood Exposure +1 (your methods smell of law‑twisting).

\section*{Act III −− Judgment on the Water}
\textbf{Goal:} Secure a settlement that prevents war, punishes the guilty, and sets the river's law.

\timer{Final Situation Timer }{8}: \textbf{Bind the Belworth Accord} \\
\timer{Counter}{6}: \textit{March at Dawn} (if Seniro's influence is high) or \textit{Honor the Fallen} (if Restorationists are ascendant)

\subsection*{Settlement Options (pick one primary + one concession)}
\begin{itemize}
  ℹtem \textbf{Hedge‑Law Trial at the River} −− Culprit faction yields its leader to a public rite at Lantern−Rites.  
    \textit{Effect:} Mandate +1; culprit Exposure fill → scandal.
  ℹtem \textbf{Indemnity Rescript} −− Black Scrip terms renegotiated; Lence gets relief; Fenwood pays in river tolls (not gold).  
    \textit{Effect:} Regency Influence +1; Syndicate gains Standing (Tag) \textit{Night‑Routes Deputized}.
  ℹtem \textbf{Red Span Purge} −− Contractors ousted; work crews handed to Syndicate and Lantern−Rites oversight.  
    \textit{Effect:} Fenwood Stability −−1 (short‑term), Mandate +1 (public cleansed).
  ℹtem \textbf{Marshal's Oath} −− Seniro swears to Constano before Sister Oria to obey lantern‑law in all crossings.  
    \textit{Effect:} House Stability +1; if broken later, automatic Exposure fill.
\end{itemize}

\subsection*{Resolution Scenes (choose 2)}
\begin{itemize}
  ℹtem \textbf{Broker the Accord} (Presence + Sway, \dv{4−−5}) −− Fill \textit{Bind the Belworth Accord} by 2 on success.
  ℹtem \textbf{Expose last‑minute sabotage} (Wits + Investigation, \dv{3−−4}) −− Prevent a counter‑timer advance.
  ℹtem \textbf{Shield the rite from blades} (Body + Melee or Stealth, \dv{3}) −− Protect the ceremony.
  ℹtem \textbf{Petition the Queen's chancellor to ratify} (Presence + Command, \dv{4}) −− Gain Mandate +1.
\end{itemize}

\subsection*{Success}
If the PCs fill \textit{Bind the Belworth Accord} before the counter‑timer triggers:
\begin{itemize}
  ℹtem \textbf{Mandate +1}
  ℹtem \textbf{Crisis −−1}
  ℹtem Apply Settlement effects.
  ℹtem Advance allied faction Influence +1 (+2 if the settlement directly satisfied their Aim).
  ℹtem \textbf{Braedon's legacy:} The old duke's Canray board is found to have a hidden compartment −− inside, a map of Aeler keystone rights that gives the Fenwoods permanent leverage on the Red Span (gain Asset: \textit{Keystone Charter}).
\end{itemize}

\subsection*{Failure}
If a counter‑timer fills first:
\begin{itemize}
  ℹtem A crossing turns bloody. \textbf{Crisis +2} (maybe +3).
  ℹtem Tick House Fenwood Stability +1 (schism risk) and Exposure +1 (butcher tales).
  ℹtem Rival faction Influence +1.
  ℹtem The Grey Wanderer appears in the aftermath, shaking his head: \textit{"He counted the exits wrong."}
\end{itemize}

\section*{Scandal & Event Prompts (use deck)}
\begin{itemize}
  ℹtem \textbf{Hearts 9} −− River Oath Broken: A ferryman violates hospitality; tick \textit{Exposure +2} on the Syndicate; start \textit{Honor Debt} [4].
  ℹtem \textbf{Spades 7} −− Night Blades: Masked men assault a witness's boat; start \textit{Witness Intimidated} [4]; on fill, your Compromise timer vanishes.
  ℹtem \textbf{Clubs 8} −− Short Rations: The Red Span crews strike; Fenwood Stability +1; tag \textit{Pickets at Dawn}.
  ℹtem \textbf{Diamonds King} −− Blood Moon on the Belworth: Public panic; all public scenes start one Position lower until appeased.
  ℹtem \textbf{Clubs Ace} −− Ghiro's ledger surfaces −− shows payments to Seniro's quartermaster; \textit{Contractors} gain \textit{Disgraced} tag; Exposure +3; Mandate +1 for exposing graft.
\end{itemize}

\section*{NPCs at a Glance}
\begin{itemize}
  ℹtem \textbf{Constano Fenwood} (heir): careful, pragmatic, favors courts. Wants a durable accord.
  ℹtem \textbf{Seniro "the Red" Fenwood} (marshal): decisive, angry, tempted by force. Wants fear to keep order.
  ℹtem \textbf{Sister Oria} (Lantern−Rites): courteous, iron ritualist. Wants peace kept in the rites.
  ℹtem \textbf{Lady−Envoy Mael Vesht} (Regency): smiling lawyer. Wants leverage for Lence's successor.
  ℹtem \textbf{Captain Pell Rivermark} (Syndicate): superstitious pilot. Wants exemptions and no riots.
  ℹtem \textbf{Master Mason Ghiro} (Contractors): perfumed operator. Wants the Red Span forever unfinished.
  ℹtem \textbf{Jaska the Poet} (Restorationists): martyr‑maker. Wants spectacle and justice.
  ℹtem \textbf{Duke Braedon Fenwood} (deceased): speaks through his Canray board and a sealed letter. The party may find his final instructions.
\end{itemize}

\subsection*{Voice & Levers (use in play):}
\begin{itemize}
  ℹtem Constano: \textit{"By hedge‑law precedent..."} (Strings: Moot Court access)
  ℹtem Seniro: \textit{"River justice is steel."} (Strings: Border patrols)
  ℹtem Oria: \textit{"Within the lantern's ring, all are guests."} (Strings: Hospitality sanctum)
  ℹtem Ghiro: \textit{"Contracts are like bridges −− they only need to hold long enough to collect."}
  ℹtem Jaska: \textit{"The dead have voices. We will lend them ours."}
\end{itemize}

\section*{Rewards, Standing & Fallout}
On a public, clean win: convert 3 favors from Lantern−Rites or the Syndicate into \textbf{Standing (Tag)}. Grant Audience Tag: \textit{Steady River}. Increase \textbf{Mandate +1} (now at least 3). Braedon's slate fragment −− Minor Asset granting +1 die to any roll involving strategy or timing, once per session.

On messy wins: rival Exposure +1 but allied Stability −−1. The settlement holds, but someone whispers: \textit{"The Fenwoods won by cheating."}

On failure: the volatile \textit{Black Scrip} stake triggers −− audits, seizures, or arrests. The river stays closed for a season. Set up the next arc.

\section*{Running Notes}
\begin{itemize}
  ℹtem \textbf{Stay Visible:} Every spend changes the map: a tag, a timer, a dial.
  ℹtem \textbf{Escalate Honestly:} When Crisis climbs, show it: protests, inquisitors, budgets, bans.
  ℹtem \textbf{Reward Face−Saving:} In intrigue, dignity can be worth more than coin.
  ℹtem \textbf{Keep It Human:} Even in a world of rites and spirits, politics is people.
  ℹtem \textbf{Use Strings as positional levers and as victory spoils.}
  ℹtem \textbf{Braedon's Canray board} can be used as a prop. The GM may reveal that the positions of pieces on his board mirror the political landscape −− the last unfinished game of a dying duke.
\end{itemize}

\section*{Quick Prep Checklist}
\begin{itemize}
  ℹtem Print Faction Sheets & Strings (update with Constano, Seniro, and Braedon's death).
  ℹtem Mark Mandate 2 / Crisis 2 (or 3).
  ℹtem Lay Act I timers: \textit{Contain the Spiral} [6]; \textit{Hotheads March} [4]; \textit{Justice Now!} [4].
  ℹtem Choose/seed 2−−3 likely culprits; plant 3 clue threads.
  ℹtem Pick 2 scandal prompts to foreshadow (e.g., Clubs 8 \textit{Short Rations}, Spades 7 \textit{Night Blades}).
  ℹtem Decide which settlement option each faction most desires.
  ℹtem Place Braedon's solar as a memory −− the party may visit to read his letter (gain a clue or a String).
\end{itemize}

\section*{Integration with \textit{The Old Duke's Wars}}
If the players have completed the vignette adventure \textit{The Old Duke's Wars}, incorporate the following echoes:
\begin{itemize}
  ℹtem If they spared the Thornhill child, the ghost haunts the Red Span, granting a one−time +1 die but adding a \textbf{Haunted} condition.
  ℹtem If they exposed Alaysha's conspiracy, Veralyn's court is more cooperative −− Mandate starts at 3 instead of 2.
  ℹtem If they executed the Thornhill child, the Restorationists are fanatical −− \textit{Justice Now!} counter‑timer starts at 2/4 instead of 0.
  ℹtem If they kept Baldvyn's letter unopened, the letter contains a secret Aeler clause −− the party gains the Keystone Charter asset immediately.
\end{itemize}
The GM may also allow players to use Braedon's Canray proverbs as a minor resource: spend 1 Boon to ask for a "Kon"reh insight" −− the GM gives a cryptic but actionable clue.

\vfill
\begin{center}\small \emph{This file presents procedures and a complete adventure. Defer to the Fate's Edge core SRD for roll math and basic adjudication.}\end{center}
\end{document}
